\section*{Hoofdstuk \ref{ch:g7:triangles} --- Driehoeken en hoeken}

\begin{solution}{exo:g7:triangles:1}
Complementen: $60^\circ$; $45^\circ$; $18^\circ$.
Supplementen: $150^\circ$; $135^\circ$; $108^\circ$.
\end{solution}

\begin{solution}{exo:g7:triangles:2}
De overstaande hoek meet ook $35^\circ$
(\cref{thm:g7:triangles:equalities}); de twee overige hoeken zijn er supplementair
mee: elk $180 - 35 = 145^\circ$.
\end{solution}

\begin{solution}{exo:g7:triangles:3}
(a) $180 - 100 = 80^\circ$. \quad
(b) $180 - 118 = 62^\circ$. \quad
(c) $180 - 150 = 30^\circ$.
\end{solution}

\begin{solution}{exo:g7:triangles:4}
Tophoek $40^\circ$: de basishoeken delen $180 - 40 = 140^\circ$ gelijk: elk
$70^\circ$.

Basishoek $70^\circ$: de twee basishoeken zijn gelijk, dus is de tophoek
$180 - 2 \times 70 = 40^\circ$.
\end{solution}

\begin{solution}{exo:g7:triangles:5}
$(6, 8, 12)$: $12 < 6 + 8 = 14$: de driehoek bestaat.

$(5, 5, 11)$: $11 > 5 + 5 = 10$: onmogelijk.

$(4, 9, 13)$: $13 = 4 + 9$: een platte ``driehoek'' --- de punten liggen op één
\omterm{def:g6:lines:objects}{lijn}, dus geen echte driehoek.

$(7, 7, 7)$: $7 < 14$: een \omterm{def:g6:shapes:triangles}{gelijkzijdige driehoek}.
\end{solution}

\begin{solution}{exo:g7:triangles:6}
$x + 2x + 3x = 6x = 180^\circ$, dus $x = 30^\circ$: de hoeken zijn $30^\circ$,
$60^\circ$ en $90^\circ$ --- een \omterm{def:g6:shapes:triangles}{rechthoekige driehoek}.
\end{solution}

\begin{solution}{exo:g7:triangles:7}
Voorspelling: $\widehat{BAC} = 180 - (50 + 60) = 70^\circ$; de gradenboog op de
afgewerkte figuur bevestigt het.
\end{solution}

\begin{solution}{exo:g7:triangles:8}
Bij het eerste snijpunt zijn de vier hoeken $54^\circ$, $126^\circ$, $54^\circ$ en
$126^\circ$ (paren \omterm{def:g7:triangles:pairs}{overstaande hoeken} en supplementen). De parallelle \omterm{def:g6:lines:objects}{lijn} herhaalt
dezelfde vier maten bij het tweede snijpunt (overeenkomstige hoeken): acht hoeken in
totaal, vier van $54^\circ$ en vier van $126^\circ$.
\end{solution}

\begin{solution}{exo:g7:triangles:9}
Zij $c$ de derde zijde. De driehoeksongelijkheid eist $c < 8 + 3 = 11$ en ook
$8 < c + 3$, dus $c > 5$. Bijgevolg $5 < c < 11$: de lengte $7$ cm werkt; de lengte
$12$ cm (of $4$ cm) niet.
\end{solution}

\begin{solution}{exo:g7:triangles:10}
$\widehat A + \widehat B = 180 - 90 = 90^\circ$, en $\widehat A = 2\widehat B$, dus
$3\widehat B = 90^\circ$: $\widehat B = 30^\circ$ en $\widehat A = 60^\circ$.
\end{solution}

\begin{solution}{exo:g7:triangles:11}
De diagonaal $[AC]$ verdeelt $ABCD$ in de driehoeken $ABC$ en $ACD$. De vier hoeken
van de \omterm{def:g4:geometry:polygon}{vierhoek} zijn precies de zes hoeken van de twee driehoeken, anders gegroepeerd
(de hoeken in $A$ en in $C$ zijn elk in twee gesplitst). Totaal:
$180 + 180 = 360^\circ$. Een \omterm{def:g4:geometry:polygon}{vijfhoek} valt vanuit één \omterm{def:g5:solids:def}{hoekpunt} in drie driehoeken
uiteen: $3 \times 180 = 540^\circ$.
\end{solution}

\begin{solution}{pb:g7:triangles:1}
\textbf{1.} De diagonalen uit één \omterm{def:g5:solids:def}{hoekpunt} van een \omterm{def:g4:geometry:polygon}{zeshoek} verdelen hem in $4$
driehoeken, dus zijn de binnenhoeken samen $4 \times 180 = 720^\circ$.

\textbf{2.} Uit één \omterm{def:g5:solids:def}{hoekpunt} gaan diagonalen naar alle \omterm{def:g5:solids:def}{hoekpunten} behalve zichzelf
en zijn twee buren: $n - 3$ diagonalen, die de \omterm{def:g4:geometry:polygon}{veelhoek} in $n - 2$ driehoeken
verdelen. Elke binnenhoek van de \omterm{def:g4:geometry:polygon}{veelhoek} wordt over de driehoeken verdeeld, dus is
de \omterm{def:g1:addition:def}{som}
\[
(n - 2) \times 180^\circ .
\]
Controle: $n = 3$: $180^\circ$; $n = 4$: $360^\circ$
(\cref{exo:g7:triangles:11}); $n = 5$: $540^\circ$; $n = 6$: $720^\circ$.

\textbf{3.} Tienhoek: $8 \times 180 = 1\,440^\circ$. Twaalfhoek:
$10 \times 180 = 1\,800^\circ$.

\textbf{4.} Door $n$ delen: driehoek $60^\circ$; vierkant $90^\circ$; \omterm{def:g4:geometry:polygon}{vijfhoek}
$\frac{540}{5} = 108^\circ$; \omterm{def:g4:geometry:polygon}{zeshoek} $\frac{720}{6} = 120^\circ$; \omterm{def:g4:geometry:polygon}{achthoek}
$\frac{1080}{8} = 135^\circ$.

\textbf{5.} Een \omterm{def:g5:solids:def}{hoekpunt} erbij voegt één driehoek aan de bundel toe: de nieuwe
\omterm{def:g4:geometry:polygon}{veelhoek} heeft $n - 1$ driehoeken nodig waar de oude er $n - 2$ nodig had. Eén
driehoek erbij, dus $180^\circ$ erbij.

\textbf{6.} De wandelaar draait over het supplement: $180^\circ - \widehat A$. Voor
de driehoek $40$--$60$--$80$ zijn de draaiingen $140^\circ$, $120^\circ$ en
$100^\circ$.

\textbf{7.} $140 + 120 + 100 = 360^\circ$: één volle draai.

\textbf{8.} Over de hele wandeling begint en eindigt je neus in dezelfde richting,
nadat hij precies één keer is rondgegaan: al het draaien in de hoeken samen is één
volledige omwenteling, $360^\circ$ --- ongeacht het aantal hoeken en de vorm van het
blok. (Het feit is prachtig onafhankelijk van $n$.)

\textbf{9.} In elke hoek is binnenhoek $+$ draaiing $= 180^\circ$. Optellen over de
$n$ hoeken:
\[
\text{(som van de binnenhoeken)} + 360^\circ = 180^\circ \times n,
\qquad\text{dus}\qquad
\text{som van de binnenhoeken} = (n - 2) \times 180^\circ
\]
--- de formule van vraag 2, opnieuw bewezen, al wandelend.

\textbf{10.} Een binnenhoek van $150^\circ$ betekent draaiingen van $30^\circ$, en
$n = \frac{360}{30} = 12$: de regelmatige twaalfhoek. Een binnenhoek van $155^\circ$
zou draaiingen van $25^\circ$ betekenen, en $\frac{360}{25} = 14.4$ is geen heel
aantal hoeken: zo'n regelmatige \omterm{def:g4:geometry:polygon}{veelhoek} bestaat niet.

\textbf{11.} Driehoek ($60^\circ$): er komen $6$ samen,
$6 \times 60 = 360$. Vierkant ($90^\circ$): $4$. \omterm{def:g4:geometry:polygon}{Zeshoek} ($120^\circ$): $3$, met
$3 \times 120 = 360$. Alle drie sluiten perfect.

\textbf{12.} Drie hoeken van een \omterm{def:g4:geometry:polygon}{vijfhoek}: $3 \times 108 = 324^\circ$ --- er blijft
een gat van $36^\circ$. Vier hoeken: $4 \times 108 = 432^\circ > 360^\circ$ ---
overlap. Geen van beide werkt: regelmatige \omterm{def:g4:geometry:polygon}{vijfhoeken} kunnen het vlak niet
betegelen.

\textbf{13.} In een betegelingshoek komen minstens $3$ tegels samen. Een regelmatige
\omterm{def:g4:geometry:polygon}{veelhoek} met meer dan zes zijden heeft binnenhoeken \emph{groter} dan $120^\circ$
(vraag 4 en het groeien van de hoek met $n$), dus zijn drie hoeken meer dan
$3 \times 120 = 360^\circ$: overlap, onmogelijk. Bij precies zes zijden werkt
$3 \times 120 = 360$; onder de zes regelen de vragen 11 en 12 de gevallen. De
volledige lijst regelmatige betegelaars: driehoek, vierkant, \omterm{def:g4:geometry:polygon}{zeshoek}.

\textbf{14.} \omterm{def:g4:geometry:polygon}{Achthoeken} en vierkanten: $135 + 135 + 90 = 360^\circ$: de hoek sluit
--- het klassieke straatpatroon. Vierkant, \omterm{def:g4:geometry:polygon}{zeshoek}, twaalfhoek: de draaiing van de
twaalfhoek is $\frac{360}{12} = 30^\circ$, dus is haar binnenhoek $150^\circ$, en
$90 + 120 + 150 = 360^\circ$: dat werkt ook.

\textbf{15.} De hoeken van de honingraat: $3 \times 120 = 360^\circ$, perfect. Van
de tegels driehoek, vierkant en \omterm{def:g4:geometry:polygon}{zeshoek} met dezelfde \omterm{def:g6:measure:area}{oppervlakte} heeft de \omterm{def:g4:geometry:polygon}{zeshoek} de
kortste rand --- ze ligt van de drie het dichtst bij de cirkel van Dido
(\cref{pb:g6:measure:1}) --- dus omsluiten zeshoekige cellen dezelfde honing met de
minste was: de bijen betegelen als meetkundigen.
\end{solution}
